\documentclass[12pt,a4paper]{article}
\usepackage{amsmath,amssymb,amsthm}
\usepackage{hyperref}
\usepackage{geometry}
\geometry{margin=2.5cm}
\usepackage{enumitem}
\usepackage{booktabs}

\newtheorem{theorem}{Theorem}[section]
\newtheorem{lemma}[theorem]{Lemma}
\newtheorem{corollary}[theorem]{Corollary}
\newtheorem{proposition}[theorem]{Proposition}
\theoremstyle{definition}
\newtheorem{definition}[theorem]{Definition}
\newtheorem{axiom_rs}{Axiom}
\theoremstyle{remark}
\newtheorem{remark}[theorem]{Remark}

\title{The Past Theorem:\\
Low-Entropy Initial Conditions as a Mathematical Necessity}

\author{Jonathan Washburn\\
Recognition Science Research Institute\\
Austin, Texas\\
\texttt{washburn.jonathan@gmail.com}}

\date{March 2026}

\begin{document}
\maketitle

\begin{abstract}
The Past Hypothesis---the postulate that the universe began in a state
of extraordinarily low entropy---is widely regarded as the most puzzling
boundary condition in cosmology.  Penrose estimated its improbability at
$1:10^{10^{123}}$.  We show that within Recognition Science~\cite{RS_Axioms},
the puzzle is structurally dissolved.  The cost functional
$J(x) = \tfrac{1}{2}(x + x^{-1}) - 1$ possesses a \emph{unique} global
minimum at $x = 1$ with $J(1) = 0$.  For any ledger configuration of
$N$ entries, the total defect $D_{\rm tot} = \sum_{i=1}^N J(x_i)$ is
zero if and only if every entry satisfies $x_i = 1$ (the unity
configuration).  This establishes: (i)~the ground state of any
cost-governed ledger is unique; (ii)~there is exactly one zero-entropy
configuration; (iii)~cost minimization selects this unique configuration
as the preferred initial state, leaving no room for alternative choices.
The proof requires no statistical mechanics, no measure theory, and no
fine-tuning.  It follows from three axioms that uniquely determine~$J$.
We carefully distinguish the RS ground state (the unique equilibrium)
from the cosmological initial state, and explain how the Past Hypothesis
is resolved by uniqueness rather than by assertion.
\end{abstract}

%% ============================================================
\section{Introduction}
\label{sec:intro}
%% ============================================================

\subsection{The Problem}

One of the deepest puzzles in the foundations of physics is the
low-entropy initial condition of the universe.  The second law of
thermodynamics tells us that entropy increases toward the future.  Running
this logic backward, the universe must have begun in a state of
extraordinarily low entropy.  But why?

Penrose~\cite{Penrose1989} quantified the specialness of the initial state
as follows.  The phase-space volume of the observable universe's current
macrostate is some $V_{\mathrm{now}}$.  The phase-space volume of the
initial state is $V_{\mathrm{init}} \approx e^{S_{\mathrm{init}}}$, where
$S_{\mathrm{init}}$ is minuscule compared to the maximum entropy
$S_{\mathrm{max}} \sim 10^{123}$.  The probability of the initial
condition, measured in phase space, is approximately
\[
  P \;\sim\; \frac{V_{\mathrm{init}}}{V_{\mathrm{max}}}
  \;\sim\; \frac{1}{10^{10^{123}}}.
\]
This number is so small that no natural process could have ``selected'' it
from an ensemble.  The standard response is simply to postulate the initial
condition---the Past Hypothesis~\cite{Albert2000}; see
also~\cite{Carroll2010,Earman2006,Callender2004} for critical
assessments from different perspectives.

\subsection{The Resolution}

We show that the cost axioms of Recognition Science force the ground
state to be a specific, unique configuration.  The key steps are:

\begin{enumerate}
  \item The cost functional $J(x) = \tfrac{1}{2}(x + x^{-1}) - 1$ is
  non-negative, vanishing only at $x = 1$.
  \item For a configuration of $N$ ledger entries, the total defect
  $D_{\mathrm{tot}} = \sum_{i=1}^{N} J(x_i)$ is the sum of non-negative
  terms.
  \item $D_{\mathrm{tot}} = 0$ if and only if every entry equals~1.
  \item The unity configuration $(1, \ldots, 1)$ is thus the unique
  global minimizer and the only zero-defect state.
\end{enumerate}

The Past Hypothesis puzzle is resolved by uniqueness: there is exactly
\emph{one} configuration of zero entropy, so the ``choice'' of initial
condition is not a free parameter but a mathematical necessity forced by
the cost axioms.  We distinguish two related but distinct claims:
\begin{itemize}[nosep]
  \item \textbf{Uniqueness Theorem (proved here)}: There is exactly one
  zero-defect configuration: the unity configuration.
  \item \textbf{Initial Condition Claim (structural, not proved here)}:
  The first recognition events initialize the ledger at the
  minimum-cost configuration consistent with their constraints, which is
  the unity configuration.
\end{itemize}
The former is mathematical; the latter is a physical application of cost
minimization.  Together, they convert the Past Hypothesis from an
arbitrary boundary condition into a uniquely selected state.

%% ============================================================
\section{The Cost Functional}
\label{sec:cost}
%% ============================================================

\begin{axiom_rs}[A1 --- Normalization]
$J(1) = 0$.
\end{axiom_rs}

\begin{axiom_rs}[A2 --- Recognition Composition Law]
$J(xy) + J(x/y) = 2J(x)J(y) + 2J(x) + 2J(y)$ for all $x, y > 0$.
\end{axiom_rs}

\begin{axiom_rs}[A3 --- Calibration]
$J''_{\log}(0) = 1$, where $J_{\log}(t) := J(e^t)$.
\end{axiom_rs}

\begin{theorem}[Cost Uniqueness]
\label{thm:Junique}
The unique continuous solution to A1--A3 is
\begin{equation}
  J(x) = \frac{1}{2}\!\left(x + \frac{1}{x}\right) - 1.
  \label{eq:J}
\end{equation}
\end{theorem}

\begin{proof}
Setting $H(t) = 1 + J(e^t)$, A2 becomes $H(s+t) + H(s-t) = 2H(s)H(t)$
(d'Alembert's equation).  A1 gives $H(0) = 1$; A3 gives $H''(0) = 1$.
By the Acz\'el classification theorem~\cite{Aczel1966}, the unique
continuous even solution is $H = \cosh$, yielding~\eqref{eq:J}.
\end{proof}

\begin{lemma}[Fundamental Properties of $J$]
\label{lem:Jprops}
For all $x > 0$:
\begin{enumerate}[label=(\roman*)]
  \item $J(x) \geq 0$ (non-negativity).
  \item $J(x) = 0 \iff x = 1$ (unique zero).
  \item $J(x) = J(1/x)$ (reciprocal symmetry).
  \item $J(x) > 0$ whenever $x \neq 1$ (strict positivity away from unity).
\end{enumerate}
\end{lemma}

\begin{proof}
(i)--(ii): $x + x^{-1} \geq 2$ by AM--GM, with equality iff $x = 1$.
(iii): $x^{-1} + x = x + x^{-1}$.  (iv): Immediate from (i) and (ii).
\end{proof}

%% ============================================================
\section{Configurations and Total Defect}
\label{sec:config}
%% ============================================================

\begin{definition}[Configuration]
A \emph{configuration} of size $N \geq 1$ is a tuple
$(x_1, x_2, \ldots, x_N) \in \mathbb{R}_{>0}^N$.
\end{definition}

\begin{definition}[Total Defect]
The \emph{total defect} (or entropy) of a configuration is
\begin{equation}
  D_{\mathrm{tot}}(x_1, \ldots, x_N) := \sum_{i=1}^{N} J(x_i)
  = \sum_{i=1}^{N} \left[\frac{1}{2}\!\left(x_i + \frac{1}{x_i}\right)
  - 1\right].
  \label{eq:Dtot}
\end{equation}
\end{definition}

\begin{lemma}[Total Defect is Non-negative]
\label{lem:Dnonneg}
$D_{\mathrm{tot}} \geq 0$ for every configuration.
\end{lemma}

\begin{proof}
Each summand $J(x_i) \geq 0$ by Lemma~\ref{lem:Jprops}(i).
\end{proof}

\begin{definition}[Unity Configuration]
The \emph{unity configuration} of size $N$ is the vector
$\mathbf{1} := (1, 1, \ldots, 1)$ in $\mathbb{R}_{>0}^N$.
\end{definition}

\begin{theorem}[Unity Has Zero Defect]
\label{thm:unity-zero}
$D_{\mathrm{tot}}(\mathbf{1}) = 0$.
\end{theorem}

\begin{proof}
$J(1) = 0$, so every summand vanishes.
\end{proof}

%% ============================================================
\section{The Past Theorem}
\label{sec:past}
%% ============================================================

\begin{theorem}[Zero Defect Characterization]
\label{thm:zero-char}
$D_{\mathrm{tot}}(x_1, \ldots, x_N) = 0$ if and only if $x_i = 1$ for
every $i \in \{1, \ldots, N\}$.
\end{theorem}

\begin{proof}
$(\Leftarrow)$: Theorem~\ref{thm:unity-zero}.

$(\Rightarrow)$: Suppose $\sum_{i} J(x_i) = 0$.  Since each $J(x_i) \geq 0$,
we must have $J(x_i) = 0$ for every~$i$.  (Otherwise, some $J(x_j) > 0$
and the sum would be positive: $\sum_i J(x_i) \geq J(x_j) > 0$.)
By Lemma~\ref{lem:Jprops}(ii), $J(x_i) = 0$ implies $x_i = 1$.
\end{proof}

\begin{theorem}[Global Minimum]
\label{thm:global-min}
For every configuration $\mathbf{x}$,
$D_{\mathrm{tot}}(\mathbf{1}) \leq D_{\mathrm{tot}}(\mathbf{x})$.
\end{theorem}

\begin{proof}
$D_{\mathrm{tot}}(\mathbf{1}) = 0 \leq D_{\mathrm{tot}}(\mathbf{x})$
by Theorem~\ref{thm:unity-zero} and Lemma~\ref{lem:Dnonneg}.
\end{proof}

\begin{theorem}[Unique Global Minimizer]
\label{thm:unique-min}
If $D_{\mathrm{tot}}(\mathbf{x}) = D_{\mathrm{tot}}(\mathbf{1})$, then
$\mathbf{x} = \mathbf{1}$.  The unity configuration is the unique
global minimizer.
\end{theorem}

\begin{proof}
$D_{\mathrm{tot}}(\mathbf{x}) = 0$ by hypothesis, so $x_i = 1$ for
all~$i$ by Theorem~\ref{thm:zero-char}.
\end{proof}

\begin{theorem}[Positive Entropy Away from Unity]
\label{thm:pos-entropy}
If $x_j \neq 1$ for some $j$, then $D_{\mathrm{tot}} > 0$.
\end{theorem}

\begin{proof}
$J(x_j) > 0$ by Lemma~\ref{lem:Jprops}(iv), and all other terms are
non-negative, so $D_{\mathrm{tot}} \geq J(x_j) > 0$.
\end{proof}

\begin{theorem}[The Past Theorem --- Ground State Uniqueness]
\label{thm:past-theorem}
The following hold simultaneously for every $N \geq 1$:
\begin{enumerate}[label=(\roman*)]
  \item There exists a unique configuration with zero total defect: the
  unity configuration $\mathbf{1} = (1,\ldots,1)$.
  \item $\mathbf{1}$ is the strict global minimizer of $D_{\mathrm{tot}}$.
  \item Every configuration $\mathbf{x} \neq \mathbf{1}$ has strictly
  positive total defect: $D_{\mathrm{tot}}(\mathbf{x}) > 0$.
\end{enumerate}
\end{theorem}

\begin{proof}
Combine Theorems~\ref{thm:zero-char}, \ref{thm:global-min},
\ref{thm:unique-min}, and~\ref{thm:pos-entropy}.
\end{proof}

\begin{corollary}[Uniqueness Resolves the Past Hypothesis]
\label{cor:past-hyp}
Because there is exactly one zero-defect configuration, the ``choice''
of low-entropy initial condition is not arbitrary.  Cost minimization
has a unique selection: the unity configuration.  The Past Hypothesis
ceases to be a puzzle requiring a fine-tuned boundary condition; it
becomes the unique output of the cost axioms.
\end{corollary}

%% ============================================================
\section{Discussion}
\label{sec:discuss}
%% ============================================================

\subsection{Why This Resolves Penrose's Puzzle}

Penrose's $1:10^{10^{123}}$ improbability assumes a measure on a set of
possible initial conditions.  In our framework, there is no such set: the
cost axioms admit exactly \emph{one} zero-cost configuration.  The
``probability'' of the initial condition is not $10^{-10^{123}}$; it is
exactly~1.  The question ``Why this initial condition rather than another?''
has no content because no alternative exists.

\subsection{Comparison with Albert's Past Hypothesis}

Albert~\cite{Albert2000} treats the Past Hypothesis as a fundamental law,
coordinate with the dynamical laws and the probability postulate.  In RS,
the Past Hypothesis is not a separate law but a consequence of the cost
axioms that already generate the dynamics.  The number of independent
postulates is reduced: where Albert needs three (dynamics + statistics +
boundary condition), RS needs only the cost axioms.

\subsection{No Fine-Tuning}

Fine-tuning requires a space of alternatives against which the actual value
is special.  The unity configuration is not ``fine-tuned'' because there is
no parameter to tune.  Every entry is 1, and no other value is admissible
at zero cost.  This is analogous to the fact that the unique real root of
$x^2 = 0$ is $x = 0$---not fine-tuned, simply forced.

\subsection{Comparison with Other Approaches}

Several approaches to the Past Hypothesis exist in the literature.
Penrose's Weyl Curvature Hypothesis~\cite{Penrose1989} constrains the
initial singularity to have vanishing Weyl tensor, enforcing low
gravitational entropy.  This is a geometric constraint analogous in spirit
to the RS approach, but it requires a separate postulate about curvature
that does not follow from general relativity.  Carroll~\cite{Carroll2010}
proposes a multiverse framework in which the low-entropy initial condition
is a typical fluctuation in a larger ensemble.  This shifts the puzzle
from ``why this initial condition?'' to ``why these multiverse dynamics?''
Callender~\cite{Callender2004} argues that the Past Hypothesis is not
genuinely puzzling because entropy is a macroscopic concept lacking
micro-physical significance.  The RS approach differs from all three: it
derives the uniqueness of the ground state from three cost axioms that
simultaneously generate the dynamics, requiring no additional postulate,
no multiverse, and no appeal to the anthropic principle.

\subsection{RS Defect, Entropy, and the Second Law}

We define the RS defect of a configuration as
$D_{\mathrm{RS}} := D_{\mathrm{tot}} = \sum_i J(x_i)$.  This quantity
measures departure from the identity ratio:
\begin{itemize}[nosep]
  \item $D_{\mathrm{RS}} \geq 0$ (non-negative).
  \item $D_{\mathrm{RS}}$ is additive over independent subsystems.
  \item $D_{\mathrm{RS}} = 0$ at the unique ground state $\mathbf{x} = \mathbf{1}$.
  \item $D_{\mathrm{RS}} > 0$ for every other configuration.
\end{itemize}

The zero-defect state $\mathbf{1}$ is the \emph{equilibrium} of the RS
ledger---the configuration of maximum symmetry and minimum cost.  In
thermodynamics, equilibrium corresponds to \emph{maximum} Boltzmann
entropy $S_B = k_B \ln \Omega$.  Therefore the two quantities are
\emph{anti-correlated}:
\[
  D_{\mathrm{RS}} = 0 \;\longleftrightarrow\; S_B = S_{B,\max}
  \quad\text{(thermodynamic equilibrium)}.
\]
As the RS ledger evolves toward the ground state, $D_{\rm RS}$ decreases
and $S_B$ increases.  The RS direction of time (decreasing $D_{\rm RS}$)
and the thermodynamic direction of time (increasing $S_B$) are the
\emph{same} direction: both point toward equilibrium.

This resolves an apparent paradox: the RS ground state has zero
\emph{defect} but maximum thermodynamic \emph{entropy}.  These are
different quantities.  The Past Hypothesis (low thermodynamic entropy
past) corresponds to the early universe having \emph{high} $D_{\rm RS}$
(many ledger entries far from unity).  The future equilibrium (maximum
thermodynamic entropy) corresponds to low $D_{\rm RS}$ (entries relaxed
to unity).  Both the RS and the thermodynamic arrow of time point from
the high-$D_{\rm RS}$ past toward the low-$D_{\rm RS}$ future.

\subsection{Clarification: Ground State vs.\ Big Bang Epoch}

\begin{remark}[Two Distinct Uses of ``Initial'']
The Past Theorem proves that the unique zero-defect state is the
RS ground state and the unique minimum-cost configuration.  This is
a logical and mathematical claim.  The physical cosmological story
is complementary:
\begin{itemize}
  \item The \emph{ground state} $\mathbf{1}$ (zero defect, formal
  minimum) is the ultimate attractor of RS dynamics and the unique
  ``initial'' configuration in the sense that cost minimization
  selects no other.
  \item The \emph{cosmological Big Bang epoch} corresponds to the first
  departure from $\mathbf{1}$: a large-defect state in which many
  ledger entries are far from unity.  This is the state of high
  $D_{\rm RS}$ (and low thermodynamic entropy) that Penrose identified.
  \item Both are consistent: the ground state is the logical starting
  point (no alternatives), and the Big Bang is the first
  \emph{physical} tick that creates high defect from that starting
  point.  The universe then evolves toward the ground state---decreasing
  $D_{\rm RS}$, increasing thermodynamic entropy---which is the second
  law.
\end{itemize}
\end{remark}

\subsection{Cosmological Implications}

The Past Theorem has several consequences:
\begin{enumerate}
  \item \textbf{Unique ground state.}  There is exactly one
  zero-defect configuration: the unity configuration.  This is the
  unique equilibrium that all RS trajectories converge toward.  The Past
  Hypothesis is explained because there is no alternative zero-defect
  state from which the universe could have ``chosen'' to start; the
  ground state is unique.

  \item \textbf{Structure corresponds to positive defect.}  Physical
  structure (particles, forces, spacetime geometry) corresponds to
  configurations with nonzero defect---departures from the featureless
  unity state.  The observed complexity of the universe is precisely the
  nonzero defect of the current configuration relative to unity.  The
  early universe (high $D_{\rm RS}$, low thermodynamic entropy) had
  more structure (departure from unity) than the future will.

  \item \textbf{Asymptotic return to unity.}  Since recognition events
  drive the ledger toward $x_i = 1$ (decreasing $D_{\rm RS}$), the far
  future of the universe is a return to the ground state.  This
  corresponds to thermodynamic maximum entropy (``heat death''): maximum
  $S_B$, minimum $D_{\rm RS}$, all ratios equilibrated to unity.

  \item \textbf{Coherent temporal picture.}  The universe's history
  runs from high $D_{\rm RS}$ (early universe, many particles and
  structures far from unity, low thermodynamic entropy) toward low
  $D_{\rm RS}$ (late universe, equilibration, high thermodynamic
  entropy).  This is the second law of thermodynamics as derived from
  RS cost minimization.
\end{enumerate}

%% ============================================================
\section{Falsifiable Predictions}
\label{sec:falsify}
%% ============================================================

\begin{enumerate}
  \item \textbf{Unique ground state.}  The Past Theorem predicts that the
  zero-defect configuration is unique: all ratios equal to unity.  If
  a physical system were found to have multiple distinct ground states
  with equal minimum cost and zero defect simultaneously, the theorem
  would be falsified.

  \item \textbf{No arbitrary initial conditions.}  RS predicts that the
  initial state of any cost-governed system is not a free parameter but
  selected by cost minimization.  A cosmological model that requires a
  continuous family of admissible initial conditions (rather than a
  unique cost-minimizing one) would be in tension with the RS framework.

  \item \textbf{Monotone defect decrease.}  For a closed system evolving
  under RS dynamics, $D_{\rm tot}$ must be non-increasing.  If an
  isolated system's RS defect were observed to increase over time
  (departing further from the unity ratio configuration), this would
  refute the cost-minimization principle.

  \item \textbf{No recurrence.}  Since the ground state $\mathbf{1}$
  is a unique global attractor and the dynamics are dissipative (not
  Hamiltonian), Poincar\'e recurrence cannot occur.  An observation of
  a fundamental-level recurrence to a higher-defect state would refute
  the framework.
\end{enumerate}

%% ============================================================
\section{Conclusion}
\label{sec:conclusion}
%% ============================================================

The Past Hypothesis---one of the most stubborn puzzles in the foundations
of physics---is structurally dissolved within Recognition Science.  The
unity configuration $(1, \ldots, 1)$ is the unique zero-cost state and
the strict global minimizer of total defect.  Penrose's estimate of
$1:10^{10^{123}}$ implicitly assumes a set of possible initial conditions
from which one is selected; in RS, no such set exists.  The cost axioms
admit exactly one zero-cost configuration, making the ``probability'' of
the initial condition exactly~1.

The resolution rests on two results: (i)~the Past Theorem (uniqueness
of the ground state), which is a mathematical consequence of the cost
axioms; and (ii)~the identification of RS defect $D_{\rm RS}$ with
thermodynamic entropy $S_B$ as anti-correlated quantities, so that
the direction of defect decrease (toward the unique ground state) is
the same as the direction of entropy increase (the second law).
Together, these convert the Past Hypothesis from an unexplained
boundary condition into a structural inevitability.

%% ============================================================
\begin{thebibliography}{99}

\bibitem{RS_Axioms}
J.~Washburn, ``The Algebra of Reality: A Recognition Science Derivation
of Physical Law,'' \emph{Axioms} \textbf{15}(2), 90 (2026).

\bibitem{Aczel1966}
J.~Acz\'el, \textit{Lectures on Functional Equations and Their Applications},
Academic Press, New York (1966).

\bibitem{Albert2000}
D.~Z.~Albert, \emph{Time and Chance} (Harvard University Press, 2000).

\bibitem{Carroll2010}
S.~Carroll, \emph{From Eternity to Here} (Dutton, 2010).

\bibitem{Penrose1989}
R.~Penrose, \emph{The Emperor's New Mind} (Oxford University Press, 1989).

\bibitem{Earman2006}
J.~Earman, ``The `Past Hypothesis': Not even false,''
\emph{Studies in History and Philosophy of Modern Physics}
\textbf{37}, 399 (2006).

\bibitem{Callender2004}
C.~Callender, ``There is no puzzle about the low-entropy past,'' in
\emph{Contemporary Debates in Philosophy of Science}, ed.\ C.~Hitchcock
(Blackwell, 2004).

\end{thebibliography}

\end{document}
